\chapter{Marco conceptual}

\section{Revisión de la literatura}

Esta sección reúne la literatura relevante para el tema de estudio. Según los lineamientos, se recomienda priorizar una base documental compuesta aproximadamente por 60\% de libros y 40\% de páginas web académicas o institucionales.

\section{Conceptos, teorías y modelos}

Aquí se desarrollan los conceptos, teorías y modelos que sustentan la monografía. Cada definición o postura teórica debe apoyarse en fuentes citadas de acuerdo con la normativa APA séptima edición.

\section{Antecedentes históricos o contextuales}

En esta sección se describen los antecedentes históricos, institucionales, tecnológicos o contextuales necesarios para comprender el problema de estudio.
